\subsection{Working capital and liquidity}

\begin{qfmethod}[The cash conversion cycle and liquidity comparison: days of receivables, inventory and payables, the operating and cash cycles, current, quick and cash ratios, and comparing liquidity across companies]\label{meth:cash-conversion-cycle}
\qffield{Statement} The cash conversion cycle (net trade cycle) measures how long cash is tied up between paying suppliers and collecting from customers \citep[sec.~19.1]{principles-of-finance-2026}:
\[ \text{CCC}=\text{days inventory}+\text{days receivables}-\text{days payables}. \]
Subramanyam and Wild define each component on consistent bases:
\begin{itemize}
\item receivables days as receivables over sales per day, compared with the credit terms;
\item inventory days as inventory over cost of sales per day;
\item payables days as payables over purchases or cost of sales per day.
\end{itemize}
A longer cycle needs more working capital to finance the same activity \citep[pp.~535--542]{financial-statement-analysis-2009}. The operating cycle is inventory days plus receivables days (p.~541).

The liquidity ratios measure how far current obligations are covered by assets convertible to cash: the current ratio, the quick ratio, which excludes inventory, and the cash ratio. Subramanyam and Wild add operating cash flow over current liabilities, and assess the composition and quality of current assets (p.~543). Financial flexibility adds what the ratios omit: unused borrowing capacity and lines of credit, saleable assets and discretionary spending that could be cut (p.~544).

\qffield{Notation} With annual figures: $\mathrm{DSO}=365\,\mathrm{AR}/\text{revenue}$, $\mathrm{DIO}=365\,\mathrm{Inv}/\mathrm{COGS}$, $\mathrm{DPO}=365\,\mathrm{AP}/\mathrm{COGS}$ (purchases if available). Shortening a component by $k$ days releases about $k$ times the daily denominator in cash.

\qffield{Practice} \emph{Harbour Tools, 2025} (\hyperref[meth:financial-ratio-families]{ratios}):
\begin{itemize}
\item DSO 58.30 days, DIO 107.88 days and DPO 91.06 days give a CCC of $58.30+107.88-91.06=75.11$ days.
\item A peer with DSO 45, DIO 70 and DPO 60 has a 55-day cycle.
\item The gap is mostly inventory. Ten fewer inventory days on Harbour's 2025 cost of sales of about 697 would release $10\times697/365=19.1$ million of cash.
\item Liquidity: current ratio 1.84, quick ratio 0.93 and cash ratio 0.16.
\end{itemize}
Harbour's longer cycle explains why its current ratio looks comfortable while its cash ratio is thin: much of its current assets are inventory, the least liquid of them \citep[p.~539]{financial-statement-analysis-2009}.

OpenStax's worked company has inventory of 36.5 days and receivables of 54.7 days, less 34 days of payables, a 57.2-day cycle \citep[sec.~19.1]{principles-of-finance-2026}.

The comparison is meaningful only within an industry. A grocer turns inventory in days and sells for cash, so a negative cycle is normal. A machine-tool maker carries months of work in progress.

\qffield{Conditions} Year-end balances can be distorted by seasonality and by window dressing, such as paying suppliers late or factoring receivables before the reporting date. Securitised or factored receivables shorten DSO without changing the business \citep[pp.~537--539]{financial-statement-analysis-2009}; supplier-finance arrangements lengthen DPO (\hyperref[meth:cash-flow-balance-sheet-quality]{cash-flow quality}). Use average balances and the same bases across companies.

\qffield{Pitfalls} Mixing bases, for example payables over sales with inventory over cost of sales. Reading a longer DPO as pure gain: stretching suppliers beyond terms damages credit standing \citep[sec.~19.1]{principles-of-finance-2026}. Judging liquidity from the current ratio alone, without the cash cycle, the composition of current assets and unused facilities.

\qffield{See also} \hyperref[meth:working-capital-policy]{working-capital policy}; \hyperref[meth:balance-sheet-ratios]{balance-sheet ratios}; \hyperref[ex:integrated-credit-decision]{Harbour credit decision}.

\end{qfmethod}

\begin{qfmethod}[Working-capital and liquidity management: conservative, moderate and aggressive investment and financing policies, trade credit and the cost of forgone discounts, receivables and inventory policy, and sources of liquidity]\label{meth:working-capital-policy}
\qffield{Statement} \emph{Investment policy.} A company chooses how much to hold in cash, receivables and inventory for a given level of sales. A conservative policy holds more, buying safety and sales at the cost of capital tied up. An aggressive policy holds less and risks stock-outs and lost sales (Interpretation).
\begin{itemize}
\item \emph{Receivables.} Credit terms usually follow industry practice. They can be loosened to stimulate sales or tightened to cut collection risk, and are monitored with an ageing schedule and the collection period against terms \citep[sec.~19.4]{principles-of-finance-2026}.
\item \emph{Inventory.} Policy balances carrying costs against ordering costs and stock-outs, using tools such as just-in-time and the economic order quantity \citep[sec.~19.5]{principles-of-finance-2026}.
\end{itemize}

\emph{Financing policy.} Working capital has a permanent base that is always needed and a temporary part that rises with seasonal activity \citep[sec.~19.1]{principles-of-finance-2026}:
\begin{itemize}
\item a \emph{matching} (moderate) policy finances the permanent part with long-term funds and the seasonal part with short-term borrowing;
\item a \emph{conservative} policy finances even the seasonal peak long-term;
\item an \emph{aggressive} policy funds part of the permanent base short-term.
\end{itemize}
Damodaran's general rule is to match the maturity, rate type and currency of liabilities to the assets they fund. Financing long-lived assets with short-term debt can make a firm technically insolvent when its value dips, even if it is solvent in value terms \citep[pp.~411--413]{applied-corporate-finance-2014}.

\emph{Trade credit} is spontaneous financing from suppliers. With terms of $d$ discount if paid within $t_d$ days, net due in $t_n$, forgoing the discount costs \citep[sec.~19.2]{principles-of-finance-2026}
\[ \frac{d}{1-d}\cdot\frac{365}{t_n-t_d}\ \text{(simple)},\qquad \Bigl(\frac{1}{1-d}\Bigr)^{365/(t_n-t_d)}-1\ \text{(compounded)}. \]
OpenStax uses a 360-day year, which gives 36.73 percent for 2/10 net 30.

\emph{Sources of liquidity.} Primary sources are cash, operating cash flow, trade credit and committed bank lines. Secondary sources, such as asset sales, renegotiating debt or drawing down equity, signal stress and may cost more \citep[p.~544]{financial-statement-analysis-2009}.

\qffield{Practice} \emph{Original.}
\begin{itemize}
\item \emph{Trade credit.} On 2/10 net 30 with a 365-day year, forgoing the discount costs 37.24 percent simple and 44.59 percent compounded. A company with a revolver at 5.5 percent should borrow to take the discount.
\item \emph{Financing a seasonal need.} Working capital is 160 in the first half and 240 in the second; long-term funds cost 6 percent, short-term borrowing 4.5 percent and surplus cash earns 2.5 percent. Annual costs are:
\begin{itemize}
\item matching: $160\times6\%+80\times4.5\%\times\tfrac12=11.40$;
\item conservative, 240 long-term with 80 invested in the first half: $14.40-1.00=13.40$;
\item aggressive, 120 long-term and the rest short-term: $7.20+(40+120)/2\times4.5\%=10.80$.
\end{itemize}
\end{itemize}
The aggressive policy saves 0.60 a year against matching. In exchange, 40 of the permanent base and the whole seasonal peak must be refinanced every year. For Harbour, whose \hyperref[ex:integrated-credit-decision]{facility decision} turned on headroom in a downside case, that saving would not compensate for the refinancing risk.

\qffield{Conditions} The comparison assumes the short-term rate stays below the long-term rate and that short-term funding remains available. When the yield curve inverts, or credit tightens in a downturn, the aggressive policy's cost advantage reverses exactly when refinancing is hardest. Committed facilities reduce, but do not remove, rollover risk because covenants can block drawings.

\qffield{Pitfalls} Treating supplier credit as free beyond the discount date. Measuring a policy's cost without its rollover risk. Cutting inventory to improve the cash cycle without the cost of lost sales. Counting an uncommitted line or a facility subject to a covenant that would be breached in a downturn as available liquidity.

\qffield{See also} \hyperref[meth:cash-conversion-cycle]{cash conversion cycle}; \hyperref[meth:capital-project-evaluation]{working capital in project cash flows}.

\end{qfmethod}
